\section{Introducción y motivación}
\label{sec:intro}

Los ecosistemas costeros e insulares dentro de los santuarios marinos tropicales, como el Parque Nacional Natural Tayrona (PNNT) en el Caribe colombiano, enfrentan un estrés sin precedentes debido al influjo de plumas de sedimentos terrígenos, microplásticos y materia orgánica particulada \cite{bayraktarov2014physical, montoya2018seasonal}. Si bien los campos macroscópicos de forzamiento hidrodinámico ---como el sistema estacional de surgencia de La Guajira--- son ampliamente monitoreados mediante implementaciones regionales de modelos oceánicos continuos (por ejemplo, ROMS, Delft3D), persiste un cuello de botella computacional crítico en la capa límite bentónica. Los marcos tradicionales de la ingeniería ambiental suelen saltar directamente a aproximaciones continuas a macroescala, tratando la captura de partículas, la deposición y la depuración bentónica mediante términos de sumidero empíricos altamente simplificados. En consecuencia, estos macromodelos oscurecen intrínsecamente los mecanismos físicos fundamentales que gobiernan las interacciones partícula--fluido y partícula--partícula a microescala.

En la cubierta coralina y en las interfaces sedimento--agua, el transporte de fluidos local transiciona estrictamente hacia regímenes con bajos números de Reynolds ($Re \ll 1$), donde dominan las fuerzas viscosas y la lubricación hidrodinámica. Investigaciones físicas globales recientes han confirmado que las estructuras bentónicas complejas actúan como trampas altamente eficientes y dependientes de la geometría para contaminantes coloidales y particulados a través de mecanismos hidrodinámicos explícitos, incluyendo la intercepción directa en la capa límite y el estancamiento por vórtices de estela \cite{mendrik2024microplastic}. En sistemas con altas fracciones de volumen de partículas, tales como los frentes de floculación de arcillas y coloides o las matrices de arena de carbonato de alta permeabilidad características de las bahías del Tayrona, las partículas no pueden modelarse como trazadores pasivos e aislados; por el contrario, sus trayectorias dependientes del tiempo están fuertemente acopladas mediante interacciones hidrodinámicas de muchos cuerpos \cite{brady1988stokesian}.

Para cerrar esta desconexión fundamental entre la micromecánica discreta y el pronóstico ambiental a macroescala, esta investigación introduce un marco computacional multiescala riguroso. Utilizamos formulaciones de Dinámica Stokesiana (SD) de código abierto altamente optimizadas \cite{brady1988stokesian, durlofsky1987dynamic, sierou2001accelerated, banchio2003accelerated} para simular explícitamente las interacciones hidrodinámicas de muchos cuerpos y las fuerzas de lubricación que gobiernan las suspensiones coloidales dentro de geometrías complejas realistas de bajo número de Reynolds. Con el fin de escalar estas perspectivas microestructurales discretas hacia métricas utilizables por los ingenieros regionales, proponemos desarrollar resolvedores numéricos personalizados de Diferencias Finitas (FD) para evaluar las ecuaciones continuas de advección--difusión--reacción (ADR). Al acoplar matemáticamente las distribuciones discretas de seguimiento de partículas con los campos escalares continuos, escalamos la física subyacente de bajo nivel para derivar tensores de dispersión hidrodinámica efectiva ($D_{\text{eff}}$) y tasas de captura dinámica ($R$) matemáticamente rigurosos. En última instancia, este paradigma de investigación básica elimina la parametrización \emph{ad hoc} del modelado ambiental regional, proporcionando las leyes físicas precisas requeridas para anclar la conservación costera del Caribe a la mecánica de fluidos determinista.
